\documentclass[aspectratio=169,11pt]{beamer}
\usepackage[utf8]{inputenc}
\usepackage[T1]{fontenc}
\usepackage[french]{babel}
\usepackage{helvet}
\renewcommand{\familydefault}{\sfdefault}
\usepackage{hyperref}
\definecolor{bleu}{HTML}{0000FF}
\definecolor{fond}{HTML}{F4F4EF}
\setbeamercolor{background canvas}{bg=fond}
\setbeamercolor{structure}{fg=bleu}
\setbeamercolor{title}{fg=bleu}
\setbeamercolor{frametitle}{fg=bleu,bg=fond}
\setbeamercolor{normal text}{fg=black}
\setbeamertemplate{navigation symbols}{}
\setbeamertemplate{itemize item}{\textcolor{bleu}{\tiny$\blacksquare$}}
\setbeamertemplate{footline}{\hfill\footnotesize\insertframenumber\hspace{1em}\vskip6pt}
\hypersetup{colorlinks=true,urlcolor=bleu,linkcolor=bleu}
\title{CEA : Commission économique pour l'Afrique}
\subtitle{Commissions économiques régionales de l'ONU}
\author{Veille institutionnelle}
\date{2026-11-03}

\begin{document}
\begin{frame}\titlepage\end{frame}

\begin{frame}{CEA : Commission économique pour l'Afrique}
\begin{itemize}
\item Siège à Addis-Abeba, créée en 1958
\item Travaille avec l'Union africaine et la BAD sur l'Agenda 2063 et la ZLECAf
\item Economic Report on Africa
\end{itemize}
\end{frame}

\begin{frame}{CEA : Commission économique pour l'Afrique : ressources à suivre}
\small
Veille hebdomadaire.
\begin{itemize}
\item \href{https://www.uneca.org}{CEA}\\{\footnotesize Site commission, publication : hebdomadaire}
\item \href{https://www.uneca.org/publications}{Publications de la CEA}\\{\footnotesize Catalogue, publication : mensuelle}
\end{itemize}
\end{frame}

\end{document}
